\section{Conclusion and Future Work}
\label{sec:conclusion}

We studied how to make DarkDiff, a diffusion-based camera ISP for low-light raw images, cheaper
to train and run. We replaced its \sdtwo{} U-Net with three distilled, pretrained U-Nets of about
330M parameters and added illumination-aware conditioning inside every region cross-attention
layer. At an equal budget of 10 epochs on SID, all three compact, illumination-guided models
obtained better SSIM and LPIPS than the original model. They also need a third of the trainable
parameters, less than half the inference memory and 40\% less inference time. We also found that
PSNR disagreed with the visual ranking at this early stage, while SSIM and LPIPS agreed with it.

Future work:
\begin{itemize}
  \item Train to convergence, and evaluate on ELD and LRD.
  \item Run ablations that add illumination conditioning to \sdtwo{} and remove it from the compact
        backbones, to separate the two effects.
  \item Use the sRGB latent as the diffusion target when no stage-1 VAE is available, and stabilise
        the large baseline with gradient clipping and learning-rate warm-up.
  \item Distil our own compact or MobileNet-style U-Nets from \sdtwo{}, so that backbones below
        150M parameters can be used with a pretrained prior.
  \item Reduce the number of sampling steps through distillation or faster samplers, and test on
        more cameras and real-world captures.
\end{itemize}
